\documentclass[../main]{subfiles}

\begin{document}
\chapter{Theory Notes 1}
\section{HITRAN linefunction model}







\section{Writing equations in \LaTeX}
It is beautifully easy to write equations, whether they occur inline with the text  $\Psi(x,t) = A e^{i(kx -\omega t)}$ or as a separate line
\begin{equation}
\Psi(x,t) =\sum_{n=1}^\infty c_n \psi_n(x) e^{-i E_n t/\hbar}.
\end{equation}
To continue a paragraph after an equation, make sure that there's no blank line between the equation and your next line of text.
Equations that are in a separate line will be numbered, and it's easy to refer to Eq.~\ref{SE} by labelling your equations  
\begin{equation}
i\hbar\frac{\partial}{\partial t}\Psi(x,t) = \hat{H} \Psi(x,t). 
\label{SE}
\end{equation}

Are you getting tired of writing begin\{eqnarray\} and end\{eqnarray\} yet?\footnote{Notice how I made a \} by putting a backslash in front of the symbol.  The curly brace has special meaning in \LaTeX, so if you want to make one, you have to let the program know that you want the symbol.}  You can simplify your life by defining abbreviated symbols for commands:
\newcommand{\beq}{\begin{equation}}
\newcommand{\eeq}{\end{equation}}
\beq
\langle x |\Psi(t)\rangle = \Psi(x,t).
\eeq
For simplicity, I'd suggest putting all of your new commands in a separate include file.

When writing a derivation with multiple steps, it is often useful to use an equation array
\begin{eqnarray}
\int_0^{2\pi}\sin^2{x}dx & = &  \int_0^{2\pi}\frac{1}{2}\left(1-\cos{2x}\right)dx \nonumber \\
& = & \left. \frac{1}{2}\left(x - \frac{\sin{2x}}{2}\right)\right|_0^{2\pi} \nonumber \\
& = & \pi. \nonumber 
\end{eqnarray}
Note that you can suppress equation numbering, which is often desirable in derivations!


The syntax for equation arrays is similar to the syntax for arrays within equations.  For example, 
\beq
\hat{H} = \left(\begin{array}{clcr}
		\Delta & \Omega & 0\\
		\Omega^*& 0 &\kappa \\
		0& \kappa^*&-\Delta\\ 
	\end{array}\right)
\eeq

There are many more complicated things you can do with equations, but this should be enough to get started! 

\section{Tables, lists, etc... etc...}
If you want to make a table, you'll need to specify the justification for each of the columns; \TeX will automatically calculate the column widths for you.  

\begin{table}[h]
\begin{center}
\begin{tabular}{|l|l|r|l|}
\hline
Name & Age & Town & Time \\
\hline
Joe Schmoe & 23 &  Rumford & 1:04\\
\hline
Mary Q. Contrary & 44 & Yarmouth & 1:12 \\
\hline
Baxter Boulevard & 32 & Portland & 1:42 \\
\hline
Lass Plaice & 60 & Lewiston & 1:57 \\
\hline
\end{tabular}
\caption{\label{raceresults}Not too many people showed up for the first, and probably last, annual Androscoggin swimming race.}
\end{center}
\end{table}

You might also, for example, want to 

\begin{enumerate}

\item make lists

\item number your lists

\end{enumerate}

\begin{itemize}

\item use bullet points

\item find better items to put in your bullet points

\end{itemize}

\section{Figures}
\begin{figure}[htbp] % float placement: (h)ere, page (t)op, page (b)ottom, other (p)age
  \centering
  % file name: C:/Users/home/Documents/Lily/Physics template/badfigure.eps
  \includegraphics[width=2.67in,height=3.51in,keepaspectratio]{example1.png}
  \caption{This is a very bad figure.}
  \label{fig:badfigure}
\end{figure}

\section{References}
test citation \cite{chang2006}


\end{document}
